\section{Introduction}\label{introduction}

The paper's main contribution is a zero-error graph-capacity theory generated by deterministic partial views on latent tuples.

When one latent fact is represented at several modifiable locations, exact recovery of the underlying state becomes a zero-error inference problem for a deterministic multi-location encoding. The current system state either identifies one authoritative value or leaves a nontrivial ambiguity class of mutually incompatible latent states. The main development studies the resulting graph/capacity theory: partial observations on multi-fact latent tuples induce non-clique confusability graphs, exact recovery becomes graph colorability, repeated composition yields strong powers and asymptotic capacity, and a fixed Lov\'asz-$\vartheta$ upper theory supplies the corresponding asymptotic ceiling. The same framework also supports a fact-side structural question about semantic determination among coordinates. The resulting theory lies at the intersection of zero-error information theory, side-information source coding, and finite converses~\cite{shannon1956zero,korner1973graphs,lovasz1979shannon,witsenhausen1976zero,slepian1973noiseless,cover2006elements}.

\medskip
\noindent\textbf{Reading roadmap.} The paper has one main theorem arc and two consequence arcs. The main arc is the zero-error graph-capacity theory generated by partial views: finite converse, graph characterization, block and asymptotic capacity, upper theory, and equality. The consequence arcs reuse the same model to study affine fact-side structure and realizability in concrete hosts. The main invariants and their roles are summarized below.

\begingroup
\small
\setlength{\tabcolsep}{4pt}
\renewcommand{\arraystretch}{1.1}
\begin{center}
\begin{tabularx}{0.98\linewidth}{@{}>{\RaggedRight\arraybackslash}p{0.18\linewidth}>{\RaggedRight\arraybackslash}p{0.28\linewidth}>{\RaggedRight\arraybackslash}p{0.31\linewidth}>{\RaggedRight\arraybackslash}p{0.13\linewidth}@{}}
\toprule
\textbf{Layer} & \textbf{Question} & \textbf{Key invariant / result} & \textbf{Section} \\
\midrule
Finite converse & What obstruction appears once ambiguity survives? & Injectivity, confusability, counting, entropy, finite error & \S\ref{sec:converse} \\
Graph/capacity arc & How do partial views control exact recovery and scaling? & Colorability, strong powers, asymptotic capacity, $\vartheta$ upper & \S\ref{sec:graph-characterization}--\S\ref{sec:equality} \\
Fact-side structure & When is one fact semantically determined by others? & Affine representable matroid on fact indices & \S\ref{sec:affine} \\
Threshold/cost corollaries & What follows at unit rate and above it? & Unit-rate threshold and manual-update gap & \S\ref{sec:rate-corollaries} \\
Realizability / readings & When can a host attain and certify unit rate? & Causal propagation + provenance observability & \S\ref{sec:requirements}--\S\ref{sec:evaluation} \\
\bottomrule
\end{tabularx}
\end{center}
\endgroup

\medskip
\noindent\textbf{Novelty boundary.} Witsenhausen's zero-error side-information problem and related graph-based source-coding work study exact recovery once the relevant graph or side-information structure is fixed. The new step here is earlier: a deterministic multi-location partial-view model on latent tuples generates a non-clique confusability graph from admissible views. Once that graph is in hand, the coloring, strong-power, Shannon-capacity, and Lov\'asz-$\vartheta$ machinery is classical. The new results are the exact recovery, success-set, weighted-success, and equality theorems for this model-generated class. Under an affine realized-state restriction, the same framework also yields a fact-side representable-matroid theorem. The finite converse layer is the setup that reaches these structures, while the threshold, update-cost, and realizability statements are corollaries and interpretations.

The induced graph class is only partly characterized here. The model already generates non-clique graphs, includes the cluster-collapse subclass on which equality is explicit, and is closed under the block-composition operation that becomes strong product on confusability graphs. In particular, whenever view-generated confusability is transitive, the induced graph is a cluster graph, so the class contains cluster graphs arising from transitive view families. The iterated binary-square family gives a concrete infinite non-clique strong-product family inside the class. A full characterization of which finite graphs arise from admissible view families is left open.

\subsection{Problem Formulation}\label{sec:encoding-problem}

The model yields an exact-consistency analogue of zero-error resolvability for modifiable encodings. Rate is measured by the number of independently writable locations; we refer to this quantity as the \emph{independent rate} and use \emph{degrees of freedom} (DOF) only as shorthand.

\noindent\textbf{Finite converse foundation.} Let $X$ be a finite latent source, let $Y$ be the deterministic observation transcript available to a resolver, and let $T$ be a finite auxiliary tag. Exact zero-error recovery from $(Y,T)$ is possible exactly when the pair map $x \mapsto (Y(x),T(x))$ is injective on the surviving ambiguity class. The same obstruction appears as confusability, counting, conditional-entropy, decoder-output, and finite-gap formulations, and it admits a budgeted finite-error extension. In particular, exact $k$-way resolution requires at least $\log_2 k$ bits of side information.\leanmeta{\LHrng{OBS}{1}{2}, \LH{OBS5}, \LH{CIA3}, \LH{PRB47}, \LH{PRB55}, \LH{PRB68}, \LH{PRB81}, \LH{PRB83}, \LH{PRB85}, \LH{PRB87}, \LH{PRB90}}

\noindent\textbf{Threshold consequence.} For deterministic multi-location encodings, the zero-incoherence threshold equals $1$: exact zero-error recovery is possible exactly in the single-source regime. The more detailed structure begins once the finite obstruction survives and acquires nontrivial confusability structure.\leanmeta{\LHrng{COH}{1}{2}, \LH{RAT1}, \LHrng{CAP}{2}{3}}

The paper then develops a \emph{Graph Characterization} for multi-fact partial observations: the confusability graph is the structural object controlling exact consistency under partial views, exact recovery is equivalent to graph colorability, and exact finite weighted success is determined by the maximum $T$-colorable induced subgraph.\leanmeta{\LH{MFT3}, \LH{MFT9}, \LH{MFT24}}

This graph characterization then lifts to strong powers and asymptotic zero-error rates. The full $n$-block confusability graph is the $n$-fold strong power of the one-shot graph; the normalized block-rate sequence converges to a real asymptotic Shannon-capacity value equal to the supremum of the finite block rates; and that value is bounded above by both the logarithm of the complement chromatic number and a fixed Lov\'asz-$\vartheta$ upper convention. The one-shot upper object matches standard orthonormal-representation, primal-PSD, and dual-theta forms.\leanmeta{\LH{MFT35}, \LH{MFT48}, \LH{MFT55}, \LH{MFT64}, \LH{MFT69}, \LH{MFT84}}

The paper also proves an \emph{Equality Characterization}. For the original clique-fiber subclass coming from a surjective label map, asymptotic Shannon capacity and the fixed Lov\'asz-$\vartheta$ upper both collapse to $\log |\mathcal B|$, where $\mathcal B$ is the fiber-label alphabet. This equality for cluster-graph structure is classical; the new point here is that transitivity of base confusability is the exact sharpness criterion in the view-family model, while meet-witnessing and fiber coherence are stronger sufficient conditions. Under these conditions, the asymptotic Shannon capacity and the fixed Lov\'asz-$\vartheta$ upper collapse to the logarithm of the number of connected components or realized transcript fibers, depending on the structure available.\leanmeta{\LH{MFT85}, \LH{MFT91}, \LH{MFT96}, \LH{MFT102}, \LH{MFT109}}

The same framework also supports a second structural question: which fact coordinates are determined by others across the realized state family? Under an affine restriction, that coordinate-dependence question becomes span membership of coordinate functionals, so the induced fact-closure is a representable matroid on fact indices. The same realized-state model therefore carries two complementary invariants: a graph on latent states and, in the affine regime, a matroid on coordinates.\leanmeta{\LHrng{AFM}{1}{5}}

\subsection{Relation to Classical Information Theory}\label{sec:connection-it}

Three classical lines are particularly relevant.

\textbf{Zero-error information theory.} Shannon's zero-error framework and its graph-theoretic refinements describe exact communication under confusability constraints~\cite{shannon1956zero,korner1973graphs,lovasz1979shannon}. Our model translates that perspective from one-shot channels to persistent encodings under edits.

\textbf{Source coding with side information.} In Slepian--Wolf theory and Witsenhausen's zero-error side-information problem, decoder side information lowers the rate or label budget needed for correct recovery~\cite{slepian1973noiseless,witsenhausen1976zero}. Here the side information is deterministic rather than stochastic, but it plays the same operational role: exact recovery is impossible unless the available tags or observations separate all remaining alternatives.

\textbf{Finite entropy converses.} Classical Fano inequalities and related entropy bounds convert decoding difficulty into information bounds~\cite{fano1961transmission,witsenhausen1975conditional,cover2006elements}. Our setting yields a deterministic finite-source analogue in which counting, conditional entropy, decoder-output entropy, and output-gap formulations become different normal forms of the same obstruction.

\subsection{Realizability and Verification}\label{sec:realizability}

The same deterministic model also raises a realizability question: when can a concrete host system realize the rate-$1$ regime? We show that two structural properties are required:
\begin{enumerate}
\item \textbf{Causal update propagation:} derived locations must update automatically when the source changes.
\item \textbf{Provenance observability:} the system must expose enough structural information to verify which locations are authoritative and which are derived.
\end{enumerate}

\subsection{Paper Organization}\label{overview}

The main theorem arc is: a deterministic finite converse identifies the obstruction; partial views lift that obstruction to a confusability graph; block composition turns that graph into strong powers; the resulting normalized rates converge to an asymptotic Shannon-capacity value; a fixed Lov\'asz-$\vartheta$ upper theory bounds that value; and a structural transitivity criterion identifies an equality subclass. The affine and realizability theorems are consequence arcs built on the same model rather than separate primary claims.

Section~\ref{sec:foundations} gives the minimum model and threshold setup. Section~\ref{sec:converse} develops the finite converse foundation. Sections~\ref{sec:graph-characterization}--\ref{sec:equality} contain the main graph-capacity arc: graph characterization, block and asymptotic capacity, upper theory, and equality characterization. Section~\ref{sec:affine} gives the affine fact-side theorem. Section~\ref{sec:rate-corollaries} records threshold and rate-complexity corollaries, Section~\ref{sec:requirements} gives the realizability theorem, and Sections~\ref{sec:evaluation} and~\ref{sec:empirical} record representative host-level readings and supplementary case-study material. A supplementary Lean 4 artifact machine-checks the converse family, graph extension, theta upper theory, affine fact-matroid theorem, threshold chain, realizability theorem, and rate-complexity arguments~\cite{demoura2021lean4}.

\subsection{Scope}\label{sec:scope}

The scope is finite facts represented at multiple locations under admissible edits. The focus is on \emph{structural facts}, meaning facts whose encoding locations are fixed when an object, declaration, or artifact is created. This isolates the zero-error regime most clearly and captures a wide class of realizations without making the theory domain-specific.

\subsection{Contributions}\label{sec:contributions}

The main contributions are grouped into the core graph-capacity arc and the consequence arcs built on the same model.
\paragraph{Core graph-capacity contributions.}
\begin{itemize}
\tightlist
\item \textbf{A finite converse foundation} for exact consistency, presented through confusability, counting, conditional-entropy, decoder-output, and entropy-gap formulations of the same finite obstruction, together with deterministic data processing and a budgeted finite-error extension.
\item \textbf{A multi-fact confusability-graph extension} in which partial observations on latent tuples produce non-clique confusability graphs, exact recovery becomes equivalent to graph colorability, success-set exactness becomes equivalent to induced-subgraph colorability, and the exact finite weighted-success value is characterized by the maximum $T$-colorable induced subgraph.
\item \textbf{A strong-power block law and asymptotic Shannon-capacity theory} in which the $n$-block confusability graph is the $n$-fold strong power of the one-shot graph and the normalized block-rate sequence converges to a real asymptotic capacity value equal to the supremum of the finite block rates.
\item \textbf{A fixed Lov\'asz-$\vartheta$ upper theory and standard equivalent forms} in which the same asymptotic capacity is bounded above by complement-chromatic and fixed Lov\'asz-$\vartheta$ upper objects, with standard orthonormal, primal-PSD, and dual-theta forms.
\item \textbf{An equality characterization} in which the clique-fiber subclass satisfies exact asymptotic equality, and that equality exports back to the original view-family model under transitivity of confusability, with a meet-witness criterion as a weaker sufficient condition and fiber coherence as a stronger one.
\end{itemize}

\paragraph{Consequence contributions.}
\begin{itemize}
\tightlist
\item \textbf{An affine fact-matroid theorem from the same latent-state framework} in which affine realized-state families induce a representable matroid on fact indices: semantic determination is span membership, and minimal determining fact sets are exactly the bases.
\item \textbf{Threshold and realizability corollaries} showing that unit rate is the unique zero-incoherence regime, that unit rate has $O(1)$ manual update cost while higher rates incur an $\Omega(n)$ lower bound, and that causal propagation together with provenance observability characterizes realizable verifiable unit rate in concrete hosts.
\item \textbf{A supplementary machine-checked artifact} verifying the theorem chain in Lean~4.
\end{itemize}


%==============================================================================
